\documentclass[12pt]{article}
\usepackage{fullpage}
%\usepackage{setspace}
\usepackage{enumitem}
\usepackage{listings,amssymb}
\usepackage{listings}
\usepackage{graphicx,xcolor}
%\doublespacing
\usepackage[T1]{fontenc}
\usepackage{newtxtext,newtxmath}
\lstset{numbers=left, numberstyle=\tiny}
\title{CS 147: Lab 03 \\ Due Date: October 8, 2026}
\author{}
\date{\today}
\begin{document}
\maketitle
%\vspace{-0.5in}
The main objective of this problem set is to design a register file using the toolchain for the course project.  
Feel free to refer
to the Verilog cheatsheet, Verilog lecture slides linked off of canvas. 

Some advice: Before starting to write any Verilog, you should do the following:
\begin{enumerate}
\item
Break down your design into sub-modules.
\item Define interfaces between these modules.
\item Draw paper and pencil schematics for these modules (these will be handed in as scanned
\texttt{schematic.pdf} file(s)).
\item Then start writing Verilog.
\end{enumerate}

\section*{Problem 1 (30 points)}
Using Verilog, design an 8-by-16-bit register file (i.e., 8 registers, each of size 16 bits). See the Verilog
interface below. It has one write port, two read ports, three register select inputs (two for read and one for
write,) a write enable, a reset and a clock input. All register state changes occur on the rising edge of the
clock. Your basic building block must be the \textbf{D-flipflop given in the provided files}. The read ports should
be all combinational logic. Do not use tri-state logic in your design.

\textbf{Design this register file such that changing the width to 32-bit or 64-bit is straightforward} (e.g., by
using parameters).

The read and write data ports are 16 bits each. The select inputs (read and write) are 3 bits each. When the
write enable is asserted (high) the selected register will be written with the data from the write port. The
write occurs on the next rising clock edge; write data cannot flow through to a read port during the same
cycle.

There is no read enable; data from the selected registers will always appear on to the corresponding read
ports.

The reset signal is synchronous and when asserted (active high), resets all the register values to 0.
The err output should be set to 1 if any register (or other sub-module) in the register file has an error. For
now, you may assume that an error only happens in a register if the input or enable is an unknown value,
and in other sub-modules if the input is an unknown value. Otherwise, it should be set to 0.

You may find Figures B.8.7 and B.8.8 in your textbook good places to start, when thinking about how to
design your register file.

You must use a hierarchical design. Design a 16-bit register first, and then put 8 of them together with
additional logic to build the register file.

Do not make any changes to the provided \texttt{ regFile\_hier.v} file.

\section*{Problem 2 (30 points)}

In Verilog, create a register file that includes internal bypassing so that results written in one cycle can be
read during the same cycle. Do this by writing an outer "wrapper" module that instantiates your existing
(unchanged) register file module; your new module will just add the bypass logic. The list of inputs and
outputs of the outer module should be the same as that of the inner module.

Do not make any changes to the provided \texttt{regFile\_bypass\_hier.v} file.

\section*{Problem 3 (10 points)}
In this problem, you will gain some experience with Synthesis.  Follow the steps outlined in
the \texttt{synth-howto.pdf} document and get setup.
Synthesize your new register file (in the same \texttt{hw3\_2} directory).
Synthesis will create a \texttt{synth} folder which will include \texttt{rf\_bypass.syn.v, area report, timing report,} etc. Do not delete this directory - it must be included in your submission.
Make sure that:
\begin{itemize}
\item
NO cells in the \texttt{synth/cell\_report.txt} has zero area in it. If you see a zero area for a module name, then that means that module had some kind of Verilog coding error and did NOT get synthesized. You must fix this problem. Look in \texttt{synth.log}, search for that module name, and look for warnings/errors.
\item
NO cells in the \texttt{synth/reference\_report.txt} has zero total area. If it does, then some sub modules underneath have failed to synthesize. Check \texttt{cell\_report.txt}.
\end{itemize}
Copy the \texttt{synth} folder back to your laptop:
\begin{small}
\begin{verbatim}
prompt> cd assignments/hw03/hw3_2
prompt>  scp -rp <your Okta ID>@coe-ee-cad15.sjsuad.sjsu.edu:code/hw03/hw3_2/synth .
\end{verbatim}
\end{small}

\section*{Testing Instructions}

To run the testbench programs that we have distributed for this lab, use:
\begin{itemize}
\item \texttt{./run test <homework> [problem]}
\item Replace \texttt{<homework>} with the assignment name (e.g., \texttt{hw3}).
\item The \texttt{[problem]} argument is optional. If omitted, all tests for the selected homework assignment will be run.
\item For more detailed output, add \texttt{-v} (e.g., \texttt{./run test -v hw3 hw3\_1}).
\begin{itemize}
\item \textbf{Warning:} \texttt{-v} output can be extremely long in some contexts.
\end{itemize}
\end{itemize}

The testbench for this problem \texttt{(tb\_regFile\_bypass\_hier.v)} generates a random set of input
signals to your module in each cycle, and compares outputs from your module with outputs that are expected
from a perfect register file bypass implementation.

If there are no errors in your design, you will see a \textit{"TEST PASSED"} message. If the testbench failed with
a \textit{"TEST FAILED WITH xx ERRORS"} message, look for error messages like \textit{"ERRORCHECK: Read data
incorrect in cycle <cycle\_number>"} in the testbench output. Above each of these error messages you will
see the inputs to your module, your outputs and the expected outputs for that cycle which can help you
debug. As in problem 1, we may run additional tests when grading.

\section*{Verilog Rules}
Read the Verilog rules document at \texttt{assignments/tools/verilog\_rules/verilog\_rules.pdf} and adhere to the conventions.
Rules adherence will be checked when you run \texttt{./run test}.

\section*{Submission Instructions}

In addition to the Verilog, you should also turn in schematics for each of your components. The schematics
may be handwritten or computer generated but must be legible if they are handwritten. The schematic files
should be named schematic.pdf and placed in the corresponding problems subdirectory (e.g.,
\texttt{hw3\_1/schematic.pdf}). \textbf{Any solution without a corresponding schematic drawing will not be graded}.
Although the schematics may seem simple for some of these components, as your project gets bigger and
bigger, you’ll find that drawing schematics of each component and the bigger picture will make your task
much, much easier.

To build a submission for grading, run:
\begin{itemize}
\item \texttt{./run build\_submission <assignment>}
\item This will automatically run \texttt{./run test <assignment>}, then run the generated submission
through the integrated autograder.
\item A submission ZIP is created in \texttt{generated\_turnins/<assignment>/}.
\item If you run \texttt{build\_submission} multiple times for the same assignment, the generated zip files will be prepended with a number. The most recent submission should be the item with the highest number.
\item \textbf{IMPORTANT:} Running \texttt{./run build\_submission <assignment>} \textbf{DOES NOT} actively submit your assignment, it only generates the file you will submit. This file must be manually uploaded to Gradescope.
\end{itemize}

\end{document}
